%% macron-stacks.tex -- proof sheet for precomposed macrons carrying
%% combining breathings and accents, as collected in
%% https://github.com/jaycrick/ag-precomposed-macrons
%%
%% GPOS mark/mkmk anchors have to place these. The macron is already in the
%% base codepoint, so a breathing sits on the macron and an accent sits on the
%% breathing -- two levels of mark-to-mark attachment.
%%
%% One source, two engines. Send the output to build/ rather than leaving it
%% beside this file: tools/proofs holds proof SOURCES, and build/ is where
%% everything generated belongs and is already gitignored.
%%
%%     mkdir -p build/proof
%%     lualatex -output-directory=build/proof -jobname=macron-stacks-lua \
%%              tools/proofs/macron-stacks.tex
%%     xelatex  -output-directory=build/proof -jobname=macron-stacks-xe \
%%              tools/proofs/macron-stacks.tex
%%
%% LaTeX will not create the output directory, and it writes beside the current
%% directory unless told otherwise. The -jobname is what lets one source serve
%% both engines: without it the second run overwrites the first.
%%
%% It picks up the fonts from build/ whether you run it from the repository
%% root or from this directory. If the font search misbehaves, replace
%% \LibPath below with a literal path such as build/ .
%%

\documentclass[11pt,a4paper]{article}

\usepackage{fontspec}
\usepackage[margin=18mm]{geometry}
\usepackage{array}
\usepackage{booktabs}

%% The HarfBuzz comparison is LuaTeX-only, so it is switched on by engine
%% rather than by hand. fontspec's Renderer= option does not exist under XeTeX,
%% which warns "Ignored LuaTeX-only feature" and then sets the section in the
%% ordinary font -- a duplicate of the rows above it that reads like a result.
%% XeTeX does its own shaping and has no node renderer to contrast against, so
%% the section has nothing to say there. Force it off with \harfbuzzfalse after
%% this line if your luaotfload has no HarfBuzz renderer either.
\newif\ifharfbuzz
\ifdefined\directlua \harfbuzztrue \else \harfbuzzfalse \fi

%% ---------------------------------------------------------------- fonts
\makeatletter
\newif\iflocalfonts
\newcommand\lib@try[1]{%
  \iflocalfonts\else
    \IfFileExists{#1LibertinusSerif-Regular.otf}{\gdef\LibPath{#1}\localfontstrue}{}%
  \fi}
\lib@try{build/}
\lib@try{../build/}
\lib@try{../../build/}
\lib@try{./}
\makeatother

\iflocalfonts
  \setmainfont{LibertinusSerif-Regular.otf}[
    Path            = \LibPath,
    ItalicFont      = LibertinusSerif-Italic.otf,
    BoldFont        = LibertinusSerif-Bold.otf,
    Ligatures       = TeX,
    Script          = Greek,
    Language        = Greek,
  ]
  \setmonofont{LibertinusMono-Regular.otf}[Path=\LibPath]
  %% same face, but with the mark-attachment features switched off
  \newfontfamily\libnomarks{LibertinusSerif-Regular.otf}[
    Path       = \LibPath,
    Script     = Greek,
    Language   = Greek,
    RawFeature = {-mark;-mkmk},
  ]
  \ifharfbuzz
    \newfontfamily\libhb{LibertinusSerif-Regular.otf}[
      Path     = \LibPath,
      Script   = Greek,
      Language = Greek,
      Renderer = HarfBuzz,
    ]
  \fi
\else
  \GenericWarning{}{macron-stacks: no build/LibertinusSerif-Regular.otf found,
    falling back to an installed Libertinus}
  \setmainfont{Libertinus Serif}[Ligatures=TeX, Script=Greek, Language=Greek]
  \setmonofont{Libertinus Mono}
  \newfontfamily\libnomarks{Libertinus Serif}[
    Script=Greek, Language=Greek, RawFeature={-mark;-mkmk}]
  \ifharfbuzz
    \newfontfamily\libhb{Libertinus Serif}[
      Script=Greek, Language=Greek, Renderer=HarfBuzz]
  \fi
\fi

%% The size the stacks are actually judged at. Big enough that a mark sitting a
%% few units off its anchor is visible rather than inferred.
\newcommand\glyphsize{\fontsize{20}{26}\selectfont}

%% One column of the sequence tables: the glyph over its codepoints, the
%% codepoints stacked the way they are typed.
%%
%% A nested tabular rather than \shortstack, and [t] rather than the default:
%% the columns carry one, two or three codepoints, and anything that centres
%% them vertically puts the glyphs themselves on different baselines across the
%% row -- which is exactly the thing this sheet exists to look at. Top
%% alignment keeps every glyph on one line and lets the codepoints hang.
%%
%% The codepoints are not wrapped in braces either. \\ inside a group is not
%% the row separator, so bracing the second argument makes the whole stack one
%% row and the breaks inside it an error. The font switch runs to the end of
%% the cell instead, which is safe because nothing follows it.
\newcommand\cell[2]{%
  \begin{tabular}[t]{@{}c@{}}{\glyphsize\strut #1}\\[1ex]\ttfamily\tiny #2\end{tabular}}

\setlength\parindent{0pt}
\pagestyle{plain}

\begin{document}

\section*{ag-precomposed-macrons}

A comprehensive list of polytonic characters with precomposed macrons for long
vowels --- the whole of that repository's README, set in Libertinus Serif.

\bigskip
{\glyphsize ᾱ ᾱ́ ᾱ̀ ᾱ̓ ᾱ̔ ᾱ̓́ ᾱ̔́ ᾱ̓̀ ᾱ̔̀ ῑ ῑ́ ῑ̀ ῑ̓ ῑ̔ ῑ̓́ ῑ̔́ ῑ̓̀ ῑ̔̀ ῡ ῡ́ ῡ̀ ῡ̓ ῡ̔ ῡ̓́ ῡ̔́ ῡ̓̀ ῡ̔̀}

\section*{Every sequence}

Three long vowels, each with the nine breathing-and-accent combinations the
README lists. The macron comes from the base codepoint; everything after it in
the column below is a combining mark stacked on top.

\subsection*{alpha U+1FB1}
\begin{tabular}{ccccccccc}
\toprule
bare & acute & grave & smooth & rough & sm+ac & ro+ac & sm+gr & ro+gr \\
\midrule
\cell{ᾱ}{1FB1} & \cell{ᾱ́}{1FB1\\0301} & \cell{ᾱ̀}{1FB1\\0300} &
\cell{ᾱ̓}{1FB1\\0313} & \cell{ᾱ̔}{1FB1\\0314} &
\cell{ᾱ̓́}{1FB1\\0313\\0301} & \cell{ᾱ̔́}{1FB1\\0314\\0301} &
\cell{ᾱ̓̀}{1FB1\\0313\\0300} & \cell{ᾱ̔̀}{1FB1\\0314\\0300} \\
\bottomrule
\end{tabular}

\subsection*{iota U+1FD1}
\begin{tabular}{ccccccccc}
\toprule
bare & acute & grave & smooth & rough & sm+ac & ro+ac & sm+gr & ro+gr \\
\midrule
\cell{ῑ}{1FD1} & \cell{ῑ́}{1FD1\\0301} & \cell{ῑ̀}{1FD1\\0300} &
\cell{ῑ̓}{1FD1\\0313} & \cell{ῑ̔}{1FD1\\0314} &
\cell{ῑ̓́}{1FD1\\0313\\0301} & \cell{ῑ̔́}{1FD1\\0314\\0301} &
\cell{ῑ̓̀}{1FD1\\0313\\0300} & \cell{ῑ̔̀}{1FD1\\0314\\0300} \\
\bottomrule
\end{tabular}

\subsection*{upsilon U+1FE1}
\begin{tabular}{ccccccccc}
\toprule
bare & acute & grave & smooth & rough & sm+ac & ro+ac & sm+gr & ro+gr \\
\midrule
\cell{ῡ}{1FE1} & \cell{ῡ́}{1FE1\\0301} & \cell{ῡ̀}{1FE1\\0300} &
\cell{ῡ̓}{1FE1\\0313} & \cell{ῡ̔}{1FE1\\0314} &
\cell{ῡ̓́}{1FE1\\0313\\0301} & \cell{ῡ̔́}{1FE1\\0314\\0301} &
\cell{ῡ̓̀}{1FE1\\0313\\0300} & \cell{ῡ̔̀}{1FE1\\0314\\0300} \\
\bottomrule
\end{tabular}

\clearpage
\section*{Marks off, marks on}

The left column is the same text with \texttt{mark} and \texttt{mkmk} disabled,
so each combining glyph falls at its own zero-width origin. The right column is
the font doing its job. The distance between them is the anchor work.

\bigskip
\begin{tabular}{lcc}
\toprule
codepoints & marks off & marks on \\
\midrule
{\ttfamily\tiny 1FB1 0314 0301} & {\glyphsize\libnomarks\strut ᾱ̔́} & {\glyphsize\strut ᾱ̔́} \\[1.5ex]
{\ttfamily\tiny 1FD1 0313 0300} & {\glyphsize\libnomarks\strut ῑ̓̀} & {\glyphsize\strut ῑ̓̀} \\[1.5ex]
{\ttfamily\tiny 1FE1 0314 0300} & {\glyphsize\libnomarks\strut ῡ̔̀} & {\glyphsize\strut ῡ̔̀} \\[1.5ex]
{\ttfamily\tiny 1FB1 0313 0301} & {\glyphsize\libnomarks\strut ᾱ̓́} & {\glyphsize\strut ᾱ̓́} \\
\bottomrule
\end{tabular}

\ifharfbuzz
\section*{Node renderer against HarfBuzz}

LuaTeX's own shaper and HarfBuzz resolve mark attachment separately. If these
two lines differ, the font is relying on something one of them does not
implement.

\bigskip
{\fontsize{17}{40}\selectfont ᾱ ᾱ́ ᾱ̀ ᾱ̓ ᾱ̔ ᾱ̓́ ᾱ̔́ ᾱ̓̀ ᾱ̔̀ ῑ ῑ́ ῑ̀ ῑ̓ ῑ̔ ῑ̓́ ῑ̔́ ῑ̓̀ ῑ̔̀ ῡ ῡ́ ῡ̀ ῡ̓ ῡ̔ ῡ̓́ ῡ̔́ ῡ̓̀ ῡ̔̀}

{\libhb\fontsize{17}{40}\selectfont ᾱ ᾱ́ ᾱ̀ ᾱ̓ ᾱ̔ ᾱ̓́ ᾱ̔́ ᾱ̓̀ ᾱ̔̀ ῑ ῑ́ ῑ̀ ῑ̓ ῑ̔ ῑ̓́ ῑ̔́ ῑ̓̀ ῑ̔̀ ῡ ῡ́ ῡ̀ ῡ̓ ῡ̔ ῡ̓́ ῡ̔́ ῡ̓̀ ῡ̔̀}
\fi

\clearpage
\section*{At size}

The four deepest stacks, grown until any drift in the anchors shows.

\bigskip
\begin{tabular}{ll}
{\ttfamily\tiny 10pt} & {\fontsize{10}{14}\selectfont ᾱ̔́ ῑ̓̀ ῡ̔̀ ᾱ̓́} \\[1ex]
{\ttfamily\tiny 14pt} & {\fontsize{14}{18}\selectfont ᾱ̔́ ῑ̓̀ ῡ̔̀ ᾱ̓́} \\[1ex]
{\ttfamily\tiny 20pt} & {\fontsize{20}{26}\selectfont ᾱ̔́ ῑ̓̀ ῡ̔̀ ᾱ̓́} \\[1ex]
{\ttfamily\tiny 28pt} & {\fontsize{28}{36}\selectfont ᾱ̔́ ῑ̓̀ ῡ̔̀ ᾱ̓́} \\[1ex]
{\ttfamily\tiny 40pt} & {\fontsize{40}{50}\selectfont ᾱ̔́ ῑ̓̀ ῡ̔̀ ᾱ̓́} \\[1ex]
{\ttfamily\tiny 58pt} & {\fontsize{58}{70}\selectfont ᾱ̔́ ῑ̓̀ ῡ̔̀ ᾱ̓́} \\
\end{tabular}

\vfill
\iflocalfonts
Font files taken from \texttt{\LibPath} in this working tree.
\else
No local build found; an installed Libertinus was used.
\fi

\end{document}
